% !PRINTED 395
\weberpage{395}
9. \gesperrt{Bezeichnen wir also mit $\nu$ die Anzahl der verschiedenen ungeraden in $D$ aufgehenden Primzahlen, so ist nach \S\,98, 12.}
\[
\begin{aligned}
(Z, R) &= 2^{\nu}, & D &\equiv 1 \ (\mathrm{mod}\ 4), \quad D \equiv 4 \ (\mathrm{mod}\ 16),\\
(Z, R) &= 2^{\nu+1}, & D &\equiv 8, 12, 16, 24, 28 \ (\mathrm{mod}\ 32),\\
(Z, R) &= 2^{\nu+2}, & D &\equiv 0 \ (\mathrm{mod}\ 32),
\end{aligned}
\]
\gesperrt{oder allgemein}
\[
(Z, R) = 2^{\lambda},
\]
\gesperrt{wo $\lambda$ dieselbe Bedeutung hat wie in \S\,104, 3.}

Diese Zahl stimmt, wie man sieht, mit der Anzahl der Stammteiler von $D$ überein, und aus 8. folgt, daß ein absoluter Normenrest $r$ diese Eigenschaft behält, wenn $r$ um ein Vielfaches von $D$ vermehrt wird.

Daraus der Satz:

10. \gesperrt{Die absoluten Normenreste sind in einer endlichen Anzahl von arithmetischen Progressionen, die nach Vielfachen von $D$ fortschreiten, enthalten.}

Die in $D$ aufgehenden ungeraden Primzahlen und in den Fällen a) und b) von Nr. 8 die Zahl 8 oder 4 heißen die \gesperrt{charakteristischen Primzahlen und Primzahlpotenzen der Ordnung}.

\begin{center}
\S\,109. \gesperrt{Die Geschlechter der Ideale.}
\end{center}

Nach dem Vorigen können wir die Gesamtheit der rationalen Zahlen $z$, mit Ausschluß derer, die zu $D$ im Zähler oder Nenner nicht relativ prim sind, in $\mu$ Klassen (Nebengruppen) zerlegen, wobei
\weq{1}{(Z, R) = \mu}
der oben bestimmte Wert ist. Hiernach ist:
\weq{2}{Z = R + R_1 + R_2 \ldots + R_{\mu - 1}.}
Wenn eine Zahl $z$ in eine dieser Nebengruppen gehört, so gehören alle mit $z$ nach dem Modul $D$ kongruenten Zahlen in dieselbe Nebengruppe.

11. \gesperrt{Die Hauptklasse $R$ der Normenreste ist bei der Multiplikation und Division eine Gruppe.}

Gehört $r$ in die Hauptklasse $R$, so gehören zwei Zahlen $z$ und $rz$ in \gesperrt{ein und dieselbe Klasse}.

Die Gesamtheit der zu $D$ teilerfremden ganzen und gebrochenen Ideale $\mathfrak{a}$ des Körpers $\Omega$ bilden auch eine Gruppe $\overline{O}$.
